\section{Functional Description and System Design}
\subsection{System decomposition and interfaces}
The system is divided into five logical subsystems: (1) golfer ranging and following control, (2) obstacle and emergency-stop safety, (3) drive and power, (4) RFID inventory, and (5) GPS, green-data and display. The Raspberry Pi coordinates high-level logic and display; a dedicated motor-control/safety path should ensure that loss of high-level commands does not cause continued motion. Exact processor, ranging radio, motor driver and RFID-reader interfaces remain to be confirmed by component selection.
\begin{figure}[ht]\centering\begin{tikzpicture}[>=Latex,node distance=7mm,every node/.style={font=\scriptsize},b/.style={draw,align=center,minimum width=21mm,minimum height=10mm}]
\node[b] (a) {Beacon range\\input};\node[b,right=of a] (b) {Estimate\\separation};\node[b,right=of b] (c) {Compute\\speed request};\node[b,right=of c] (d) {Apply safety\\interlock};\node[b,right=of d] (e) {Drive\\motors};
\draw[->] (a)--node[above]{range}(b);\draw[->](b)--node[above]{distance}(c);\draw[->](c)--node[above]{request}(d);\draw[->](d)--node[above]{safe command}(e);
\node[b,below=of d] (s) {Obstacle +\\E-stop signals};\draw[->](s)--(d);
\end{tikzpicture}\caption{Fully decomposed example of the following-control function, including safety interlock.}\end{figure}
\subsection{Functional-to-logical allocation}
\begin{table}[ht]\centering\small\caption{Function allocation and interface labels.}\begin{tabularx}{\linewidth}{p{.25\linewidth}p{.28\linewidth}X}\toprule Function & Logical element & Interface / requirement\\\midrule Estimate target range & Ranging receiver & Beacon measurements; R-IF01\\ Compute motion request & Following controller & Distance error; R-F01\\ Inhibit unsafe motion & Safety controller & Obstacle/E-stop; R-S01--R-S03\\ Audit club slots & RFID reader / inventory service & Tag events; R-C01\\ Calculate green range & GPS/display service & Position + map data; R-G01\\\bottomrule\end{tabularx}\end{table}
\subsection{Design choices and alternatives}
Sensor-based tracking avoids camera lighting and training dependencies, but reliable target identification must be demonstrated experimentally. RFID tag placement and reader geometry will be tested against NFC or per-slot detection if bulk reads fail. GPS green-distance display is simpler than shot tracking but depends on valid green coordinates. A separate safety inhibit is preferred over relying only on the display computer; wiring and fail-safe behavior must be checked during prototype development.
\subsection{Identified prototypes}
P01 drive/power bench rig; P02 RFID bag-slot rig; P03 following/ranging rig; P04 obstacle/E-stop rig; P05 GPS/display rig; P06 integrated stationary cart; P07 moving outdoor MVP. Each has a test and milestone in Sections 5--6.\par
